% This document is compiled using pdfLaTeX
\documentclass[11pt,a4paper]{article}

% --- Chinese support ---
\usepackage[UTF8]{ctex}
% --- Compact page geometry ---
\usepackage[top=2cm, bottom=2cm, left=2cm, right=2cm]{geometry}
% --- Tables ---
\usepackage{booktabs}
\usepackage{array}
\usepackage{tabularx}
% --- Hyperlinks ---
\usepackage[hidelinks]{hyperref}
% --- Lists ---
\usepackage{enumitem}
% --- Caption ---
\usepackage[font=small,skip=4pt]{caption}

% Compact spacing
\setlength{\parskip}{2pt}
\setlength{\parindent}{0pt}
\setlist{nosep,leftmargin=*}
\renewcommand{\arraystretch}{1.15}

\title{\bfseries\large Microelectronics Journal 调研报告}
\author{}
\date{\small 2026年6月9日}

\begin{document}

\maketitle

\section{Microelectronics Journal 投稿要求}

\subsection{基本信息}

\begin{table}[!htbp]
\centering
\small
\begin{tabular}{ll}
\toprule
\textbf{项目} & \textbf{内容} \\
\midrule
期刊全称 & Microelectronics Journal (Elsevier) \\
ISSN & 0959-8324 (Print) / 1879-2391 (Online) \\
Impact Factor & 2.3 (2024) \quad CiteScore: 4.3 \\
出版周期 & 月刊（12期/年） \\
索引收录 & SCIE, Scopus \\
投稿系统 & \url{https://www.editorialmanager.com/MEJ} \\
OA费用 & USD 2,550（选择OA时） / 传统模式无版面费 \\
\bottomrule
\end{tabular}
\caption{期刊基本信息}
\end{table}

\subsection{审稿周期}
\small
初筛 $\sim$1 天，投稿至接收 $\sim$69 天（中位数），投稿至出版通常 3--6 个月。

\subsection{文章类型}
\small
Original Papers、\textbf{Review / Survey Papers}、Case Studies、Conference Reports、Book Reviews、Letters to the Editors。

\subsection{手稿格式要求}
\begin{table}[!htbp]
\centering
\small
\begin{tabular}{ll}
\toprule
\textbf{项目} & \textbf{要求} \\
\midrule
页数限制 & 无明确限制，要求简洁 \\
语言 & 英文 \\
同行评审 & Single-anonymized（单盲） \\
摘要 & 150--250 词 \\
关键词 & 最多 6 个，美式拼写 \\
参考文献 & 编号制，方括号标注 {[}1{]}, {[}2{]} \\
图表 & 分辨率 $\geq$ 300 dpi，表格须可编辑 \\
LaTeX & 推荐 elsarticle document class \\
\bottomrule
\end{tabular}
\caption{手稿格式要求}
\end{table}

\subsection{综述文章特殊要求}
\small
投稿综述前需先与编辑讨论主题和内容（\textit{``Authors must first discuss the subject and content with the Editor before submitting a review article''}）。

\section{Microelectronics Journal 近期综述文章情况}

\small
经检索，Microelectronics Journal 近年（2022--2026 年 6 月）发表的综述文章主要集中在具体器件/电路/系统级技术的深度评述，列表如下。

\begin{table}[!htbp]
\centering
\small
\renewcommand{\arraystretch}{1.3}
\begin{tabularx}{\textwidth}{c>{\raggedright\arraybackslash}Xcc}
\toprule
\textbf{序号} & \textbf{标题} & \textbf{年份} & \textbf{DOI} \\
\midrule
1 & A comprehensive review of AlGaN/GaN High electron mobility transistors: Architectures and field plate techniques for high power/high frequency applications & 2023 & 10.1016/j.mejo.2023.105951 \\
\addlinespace
2 & Advances in neuromorphic devices for the hardware implementation of neuromorphic computing systems for future artificial intelligence applications: A critical review & 2022 & 10.1016/j.mejo.2022.105634 \\
\addlinespace
3 & Revolutionizing wireless communication: A review perspective on design and optimization of RF MEMS switches & 2023 & 10.1016/j.mejo.2023.105891 \\
\addlinespace
4 & A survey on semiconductor wafer yield prediction by artificial intelligence & 2026 & 10.1016/j.mejo.2026.106958 \\
\bottomrule
\end{tabularx}
\caption{Microelectronics Journal 2022--2026 年发表的综述/评述文章}
\end{table}

\section{综述文章篇幅与结构}

\subsection{篇幅}
\small
综述的篇幅分别为：AlGaN/GaN HEMTs 综述 $\sim$15 页，Neuromorphic devices 综述 $\sim$20 页，RF MEMS switches 综述 $\sim$15 页，Wafer yield prediction 综述 $\sim$15 页（期刊最终版）。

\subsection{典型结构}
\begin{enumerate}
    \item \textbf{Introduction} —— 背景与动机
    \item \textbf{Technical Background} —— 技术原理概述
    \item \textbf{Classification / Taxonomy} —— 分类体系
    \item \textbf{Comparative Analysis} —— 各方法比较
    \item \textbf{Challenges \& Future Directions} —— 挑战与展望
    \item \textbf{Conclusion} —— 结论
    \item \textbf{References} —— 参考文献（通常 80--200 篇）
\end{enumerate}

\subsection{参考文献格式}
\small
MEJ 采用编号制，正文用方括号标注：
\begin{verbatim}
[1] T.S. Cubitt, D. Perez-Garcia, M.M. Wolf,
    Undecidability of the spectral gap,
    Nature 528 (2015) 207-211.
\end{verbatim}

\section{学术界对 Microelectronics Journal 的评价}

\subsection{核心量化指标}

\begin{table}[!htbp]
\centering
\small
\begin{tabular}{ll}
\toprule
\textbf{指标} & \textbf{数值} \\
\midrule
JCR 分区（工程：电子与电气） & Q3 \\
JCR 分区（纳米科技） & Q4 \\
中科院分区 & 工程技术 3区 \\
Scopus CiteScore 分区 & 电气工程 Q2 \\
SJR & 0.392 \\
H-index & 83 \\
自引率 & $\sim$30.4\%（偏高） \\
年发文量 & $\sim$345--363 篇 \\
LetPub 综合评分 & 5.8 / 10（声誉 6.6，影响力 4.5，速度 9.4） \\
\bottomrule
\end{tabular}
\caption{期刊核心量化指标}
\end{table}

\subsection{正面评价}

\begin{itemize}
    \item \textbf{历史悠久}：创刊于 1969 年，是微电子领域的老牌期刊，编委会阵容较强。主编为罗切斯特大学 Prof.\ E.G.\ Friedman（VLSI 设计领域知名学者），编委来自 NUS、NTU、Imperial College、Technion、UCLA 等一流机构。
    \item \textbf{录用率高}：多来源显示录用率在 85\% 以上，在国内学术社区被称为``毕业神器''、``闭眼冲''级别期刊。
    \item \textbf{审稿速度快}：平均审稿周期约 2--3 个月。编辑处理效率获多位投稿者称赞，部分案例 2 个月内完成从投稿到接收。
    \item \textbf{审稿意见专业}：多位投稿者反馈审稿人会提供``数十条建设性意见''，有效帮助提升论文质量。
    \item \textbf{对国人友好}：审稿流程正规，编辑响应迅速，适合国内作者。
    \item \textbf{不强制要求流片}：如果创新性和理论推导充分，纯仿真工作也可接受，无需流片测试。
    \item \textbf{无强制版面费}：选择传统订阅模式发表无需缴费。
\end{itemize}

\subsection{负面评价与注意事项}

\begin{itemize}
    \item \textbf{分区较低}：JCR Q3、中科院 3 区，不属于顶级期刊。
    \item \textbf{自引率偏高}（$\sim$30.4\%）：这一数值在被引分析中可能引起注意，需控制引用行为。
    \item \textbf{审稿人匹配有时较慢}：部分投稿者反馈需要发出 6--8 次审稿邀请才能找到 2--3 位审稿人，导致周期延长至 5--7 个月。
    \item \textbf{修回期限较短}：有的修回仅给 2 周时间，对于需要补充较多实验的稿件压力较大。
    \item \textbf{影响因子偏低}：IF 2.3 在该领域处于中下游水平，2024 年虽有上升趋势但整体仍偏低。
    \item \textbf{Gold OA 比例低}（$\sim$1.2\%）：开放获取文章占比极低，可能影响可见度。
\end{itemize}

\subsection{综合定位}

Microelectronics Journal 是微电子领域的一个\textbf{中档（Q3/中科院3区）但稳定可靠的 SCI 期刊}。它录用率高、审稿快、对国人友好。其影响因子近年呈上升趋势（从 2018 年的 $\sim$1.5 升至 2024 年的 2.3），表明影响力在逐步提升。

若论文工作质量较好，可考虑冲击更高分区期刊（如 IEEE Transactions on Circuits and Systems、IEEE Journal of Solid-State Circuits、IEEE Transactions on Electron Devices 等）；若目标是快速发表，MEJ 是性价比较高的选择。

\end{document}
